\documentclass[lettersize,journal]{IEEEtran}
\usepackage{amsmath,amsfonts}
\usepackage{algorithmic}
\usepackage{algorithm}
\usepackage{array}
\usepackage[caption=false,font=normalsize,labelfont=sf,textfont=sf]{subfig}
\usepackage{textcomp}
\usepackage{stfloats}
\usepackage{url}
\usepackage{verbatim}
\usepackage{graphicx}
\usepackage{cite}
\hyphenation{op-tical net-works semi-conduc-tor IEEE-Xplore}
% updated with editorial comments 8/9/2021

\begin{document}

\title{A Sample Article Using IEEEtran.cls\\ for IEEE Journals and Transactions}

\author{IEEE Publication Technology,~\IEEEmembership{Staff,~IEEE,}
        % <-this % stops a space
\thanks{This paper was produced by the IEEE Publication Technology Group. They are in Piscataway, NJ.}% <-this % stops a space
\thanks{Manuscript received April 19, 2021; revised August 16, 2021.}}

% The paper headers
\markboth{Journal of \LaTeX\ Class Files,~Vol.~14, No.~8, August~2021}%
{Shell \MakeLowercase{\textit{et al.}}: A Sample Article Using IEEEtran.cls for IEEE Journals}

\IEEEpubid{0000--0000/00\$00.00~\copyright~2021 IEEE}
% Remember, if you use this you must call \IEEEpubidadjcol in the second
% column for its text to clear the IEEEpubid mark.

\maketitle

\begin{abstract}
This document describes the most common article elements and how to use the IEEEtran class with \LaTeX \ to produce files that are suitable for submission to the IEEE.  IEEEtran can produce conference, journal, and technical note (correspondence) papers with a suitable choice of class options. 
\end{abstract}

\begin{IEEEkeywords}
Article submission, IEEE, IEEEtran, journal, \LaTeX, paper, template, typesetting.
\end{IEEEkeywords}

\section{Introduction}
% motivation
\IEEEPARstart{E}{fficient}
 motion planning in uncertain environments with safety \& efficiency guarantees


% related works; to to be classified according to our contribution
conformal prediction for safe motion planning~\cite{lindemann2023safeplanning, dixit2023adaptive,lekeufack2024decision,mei2025perceive,shin2025egocentric,contreras2025safe,sundarsingh2026safe}


% what makes functional CP beneficial in motion planning?
\begin{itemize}
    \item representation perspective (environmental geometries often admit low-dimensional representations)

    \item recent worst-case definition of score functions for motion planning are intractable, thus requiring severe approximation in practice, e.g.~\cite{mei2025perceive,sundarsingh2026safe}.
\end{itemize}

\section{Related Works}
\begin{itemize}
    \item robust control: DRO~\cite{long2026sensor}
\end{itemize}


\section{Related Works}
\begin{itemize}
    \item navigation in dynamic environments
    \begin{itemize}
        \item perception (localization, mapping, detection, tracking, $\cdots$) \& prediction in dynamic environments
        \begin{itemize}
            \item SLAM $+$ MOT~\cite{wang2007simultaneous,coue2006bayesian}

            \item particle-based~\cite{danescu2011modeling,tanzmeister2014grid}
            
            \item random finite sets~\cite{nuss2018random,chen2023continuous}; based on~\cite{mahler2007statistical}

            \item SDF learning~\cite{park2019deepsdf,ortiz2022isdf}; mostly in 3d; may be irrelevant to ours
        \end{itemize}
    \end{itemize}

    
\end{itemize}

\section{Related Works}
\begin{itemize}
    \item navigation in dynamic environments
    \begin{itemize}
        \item perception (localization, mapping, detection, tracking, $\cdots$) \& prediction in dynamic environments
        \begin{itemize}
            \item SLAM $+$ MOT~\cite{wang2007simultaneous,coue2006bayesian}

            \item particle-based~\cite{danescu2011modeling,tanzmeister2014grid}
            
            \item random finite sets~\cite{nuss2018random,chen2023continuous}; based on~\cite{mahler2007statistical}

            \item SDF learning~\cite{park2019deepsdf,ortiz2022isdf}; mostly in 3d; may be irrelevant to ours
        \end{itemize}
    \end{itemize}

    
\end{itemize}



%{\appendices
%\section*{Proof of the First Zonklar Equation}
%Appendix one text goes here.
% You can choose not to have a title for an appendix if you want by leaving the argument blank
%\section*{Proof of the Second Zonklar Equation}
%Appendix two text goes here.}

\section{Problem Formulation}




Let $\mathcal{Q}$ be the configuration space of a robot, and  $x = x(q) \in \mathcal{X} \subseteq \mathbb{R}^d$ the position in the world corresponding to $q\in \mathcal{Q}$ (i.e., the forward kinematics of the robot), where $d = 2$ or $d=3$. Let $\mathcal{O}_t \subseteq \mathcal{X}$ be the union of \emph{dynamic} obstacles (which is often represented as an occupancy map in practice).
To model a limited perception range, we follow the formulation of~\cite{janson2018safe} and introduce a perception map $P : \mathcal{X} \times \mathscr{C}(\mathcal{X}) \to \mathscr{M}_{\text{partial}}(\mathcal{X})$, where $\mathscr{C}(\mathcal{X})$ is the collection of all closed subsets of $\mathcal{X}$ (each $\mathcal{C} \in \mathscr{C}(\mathcal{X})$ being a union of obstacles, i.e.\ the \emph{unknown} occupancy of the world), and $\mathscr{M}_{\text{partial}}(\mathcal{X})$ is the set of all measurable \emph{partial} functions on $\mathcal{X}$,
\begin{equation*}
\mathscr{M}_{\text{partial}}(\mathcal{X}) \coloneqq \bigl\{ D : \mathcal{Y} \to \mathbb{R} \;\big|\; \mathcal{Y}\subseteq\mathcal{X},\ D\ \text{measurable} \bigr\}.
\end{equation*}
Intuitively, $P(x, \mathcal{C})$ is the signed distance field (SDF) sensed from configuration $x$ when the true geometry is $\mathcal{C}$; in practice it is realized as a truncated SDF~\cite{newcombe2011kinectfusion}, which we prefer to a binary occupancy map for its smooth, distance-aware structure. The use of \emph{partial} functions accounts for the finite sensing range.

Let $\mathcal{C}_t$, $t \geq 0$, denote the time-varying occupancy, whose dynamics reflect the motion of the environment. At time $t$ the robot at $x_t = x(q_t)$ receives the observation $\widehat{D}_t = P(x_t, \mathcal{C}_t)$ with domain $\mathcal{D}_t$, and the region sensed up to time $t$ is
\begin{equation*}
\mathcal{D}_{\leq t} \coloneqq \bigcup_{\tau \leq t}\mathcal{D}_{\tau}.
\end{equation*}
On $\mathcal{D}_{\leq t}$ the observations are fused into an SDF estimate $D_{\leq t}: \mathcal{X} \to \mathbb{R}$, which is set to $0$ on the unexplored region $\mathcal{X} \setminus \mathcal{D}_{\leq t}$.



Since $\mathcal{O}_{t+i}$ for $i > 0$ are inherently unobservable at time $t$, all we can do is estimate them from historical or contextual data, as modern prediction models do. Let us denote the $i$-step ahead prediction of the occupancy made at $t$ by $\mathcal{O}_{t+i|t}$ 
For safe motion planning, we are interested in correctly estimating the (signed) distance function at $t+i$:
\begin{equation*}
D_{t+i}(x) \coloneqq d(x, \mathcal{O}_{t+i}), \quad x \in \mathcal{X}
\end{equation*}
and its $i$-step ahead prediction be:
\begin{equation*}
D_{t+i|t}(x) \coloneqq d(x, \mathcal{O}_{t+i|t}), \quad x \in \mathcal{X}
\end{equation*}

\section{References Section}
You can use a bibliography generated by BibTeX as a .bbl file.
 BibTeX documentation can be easily obtained at:
 http://mirror.ctan.org/biblio/bibtex/contrib/doc/
 The IEEEtran BibTeX style support page is:
 http://www.michaelshell.org/tex/ieeetran/bibtex/
 
 % argument is your BibTeX string definitions and bibliography database(s)
%\bibliography{IEEEabrv,../bib/paper}
%
\section{Simple References}
You can manually copy in the resultant .bbl file and set second argument of $\backslash${\tt{begin}} to the number of references
 (used to reserve space for the reference number labels box).

\begin{thebibliography}{1}
\bibliographystyle{IEEEtran}

\bibitem{ref1}
{\it{Mathematics Into Type}}. American Mathematical Society. [Online]. Available: https://www.ams.org/arc/styleguide/mit-2.pdf

\bibitem{ref2}
T. W. Chaundy, P. R. Barrett and C. Batey, {\it{The Printing of Mathematics}}. London, U.K., Oxford Univ. Press, 1954.

\bibitem{ref3}
F. Mittelbach and M. Goossens, {\it{The \LaTeX Companion}}, 2nd ed. Boston, MA, USA: Pearson, 2004.

\bibitem{ref4}
G. Gr\"atzer, {\it{More Math Into LaTeX}}, New York, NY, USA: Springer, 2007.

\bibitem{ref5}M. Letourneau and J. W. Sharp, {\it{AMS-StyleGuide-online.pdf,}} American Mathematical Society, Providence, RI, USA, [Online]. Available: http://www.ams.org/arc/styleguide/index.html

\bibitem{ref6}
H. Sira-Ramirez, ``On the sliding mode control of nonlinear systems,'' \textit{Syst. Control Lett.}, vol. 19, pp. 303--312, 1992.

\bibitem{ref7}
A. Levant, ``Exact differentiation of signals with unbounded higher derivatives,''  in \textit{Proc. 45th IEEE Conf. Decis.
Control}, San Diego, CA, USA, 2006, pp. 5585--5590. DOI: 10.1109/CDC.2006.377165.

\bibitem{ref8}
M. Fliess, C. Join, and H. Sira-Ramirez, ``Non-linear estimation is easy,'' \textit{Int. J. Model., Ident. Control}, vol. 4, no. 1, pp. 12--27, 2008.

\bibitem{ref9}
R. Ortega, A. Astolfi, G. Bastin, and H. Rodriguez, ``Stabilization of food-chain systems using a port-controlled Hamiltonian description,'' in \textit{Proc. Amer. Control Conf.}, Chicago, IL, USA,
2000, pp. 2245--2249.

\end{thebibliography}


\newpage

\section{Biography Section}
If you have an EPS/PDF photo (graphicx package needed), extra braces are
 needed around the contents of the optional argument to biography to prevent
 the LaTeX parser from getting confused when it sees the complicated
 $\backslash${\tt{includegraphics}} command within an optional argument. (You can create
 your own custom macro containing the $\backslash${\tt{includegraphics}} command to make things
 simpler here.)
 
\vspace{11pt}

\bf{If you include a photo:}\vspace{-33pt}
\begin{IEEEbiography}[{\includegraphics[width=1in,height=1.25in,clip,keepaspectratio]{fig1}}]{Michael Shell}
Use $\backslash${\tt{begin\{IEEEbiography\}}} and then for the 1st argument use $\backslash${\tt{includegraphics}} to declare and link the author photo.
Use the author name as the 3rd argument followed by the biography text.
\end{IEEEbiography}

\vspace{11pt}

\bf{If you will not include a photo:}\vspace{-33pt}
\begin{IEEEbiographynophoto}{John Doe}
Use $\backslash${\tt{begin\{IEEEbiographynophoto\}}} and the author name as the argument followed by the biography text.
\end{IEEEbiographynophoto}




\vfill

\end{document}


